\chapter{\eh{world-map} Capa 3 --- Red (Network Layer)}

\section{\eh{thinking-face} ¿Qué hace la capa de red?}

\begin{concepto}
La capa de red \textbf{enruta paquetes} de una red a otra hasta llegar al
destino final, aunque el camino pase por decenas de routers.
Usa \textbf{direcciones IP} para identificar origen y destino a nivel global.

\emoji{light-bulb} \textbf{Analogía:} Las MACs son el número de asiento en tu
salón (LAN local). Las IPs son como tu dirección postal completa (funcionan
en cualquier parte del mundo). Un router es el servicio de mensajería
que sabe cómo llevar tu paquete de ciudad en ciudad hasta la dirección correcta.
\end{concepto}

\section{\eh{puzzle-piece} IPv4}

\subsection{Header IPv4 (20 bytes mínimo)}

\begin{lstlisting}
 0         1         2         3
 01234567890123456789012345678901
 ┌───────┬────┬────────────┬─────────────────────────┐
 │Ver=4  │IHL │    DSCP    │ECN│   Total Length       │
 │ 4 bits│4 b │  6 bits    │2b │     16 bits          │
 ├───────────────────────────┬───┬────────────────────┤
 │     Identification        │Flg│  Fragment Offset    │
 │        16 bits            │3b │     13 bits         │
 ├────────────┬──────────────┴───┴────────────────────┤
 │ TTL (8b)  │ Protocol (8b) │   Header Checksum      │
 ├────────────┴──────────────┴────────────────────────┤
 │               Source IP Address (32 bits)          │
 ├────────────────────────────────────────────────────┤
 │             Destination IP Address (32 bits)        │
 └────────────────────────────────────────────────────┘
\end{lstlisting}

Campos clave:
\begin{itemize}
    \item \textbf{TTL} (Time to Live): cada router decrementa en 1. Al llegar
          a 0, el paquete se descarta y se envía ICMP ``Time Exceeded'' al origen.
          Evita loops infinitos. Valor inicial típico: 64 (Linux) o 128 (Windows).
    \item \textbf{Protocol:} qué viene en el payload.
          1 = ICMP, 6 = TCP, 17 = UDP, 89 = OSPF.
    \item \textbf{Flags:} DF (Don't Fragment), MF (More Fragments).
    \item \textbf{Total Length:} tamaño total del paquete IP incluyendo header y datos.
\end{itemize}

\subsection{Clases de IP (histórico) y CIDR}

\begin{table}[h!]
\centering
\renewcommand{\arraystretch}{1.4}
\begin{tabular}{lllll}
\toprule
\textbf{Clase} & \textbf{Rango primer octeto} & \textbf{Máscara default} & \textbf{Hosts por red} & \textbf{Para} \\
\midrule
A & 1 -- 126      & /8  (255.0.0.0)     & 16,777,214 & Grandes ISPs \\
B & 128 -- 191    & /16 (255.255.0.0)   & 65,534     & Universidades \\
C & 192 -- 223    & /24 (255.255.255.0) & 254        & Empresas pequeñas \\
D & 224 -- 239    & --- (multicast)     & ---        & Streaming, routing \\
E & 240 -- 255    & --- (experimental)  & ---        & Reservado \\
\bottomrule
\end{tabular}
\caption{Clases de IPv4 (sistema clásico, reemplazado por CIDR)}
\end{table}

\begin{concepto}
\textbf{CIDR} (Classless Inter-Domain Routing, RFC 1519, 1993) reemplazó el
sistema de clases. Permite máscaras de cualquier longitud (/1 -- /32),
lo que hace un uso mucho más eficiente del espacio de IPs.

Ejemplo: \texttt{192.168.1.0/26} significa que los primeros 26 bits son la red
y los últimos 6 bits son para hosts → $2^6 - 2 = 62$ hosts por subred.
\end{concepto}

\subsection{Subnetting paso a paso}

\begin{ejemplo}
\textbf{Problema:} Dividir \texttt{192.168.10.0/24} en \textbf{4 subredes iguales}.

\textbf{Paso 1:} Necesitamos $2^n \geq 4$ subredes → $n = 2$ bits prestados.

\textbf{Paso 2:} Nueva máscara: /24 + 2 = \textbf{/26}
\[ 255.255.255.\underbrace{11}_{\text{red}}000000_2 = 255.255.255.192 \]

\textbf{Paso 3:} Tamaño de cada subred: $2^{32-26} = 2^6 = 64$ IPs por bloque.

\textbf{Paso 4:} Las 4 subredes:
\begin{itemize}
    \item \texttt{192.168.10.0/26}   → hosts \texttt{.1} a \texttt{.62}   (broadcast: \texttt{.63})
    \item \texttt{192.168.10.64/26}  → hosts \texttt{.65} a \texttt{.126}  (broadcast: \texttt{.127})
    \item \texttt{192.168.10.128/26} → hosts \texttt{.129} a \texttt{.190} (broadcast: \texttt{.191})
    \item \texttt{192.168.10.192/26} → hosts \texttt{.193} a \texttt{.254} (broadcast: \texttt{.255})
\end{itemize}

\textbf{Fórmulas:}
$\text{subredes} = 2^n$ \quad $\text{hosts/subred} = 2^h - 2$ \quad donde $n+h = 32 - \text{prefijo original}$
\end{ejemplo}

\subsection{IPs especiales (RFC 1918 y otros)}

\begin{table}[h!]
\centering
\renewcommand{\arraystretch}{1.3}
\begin{tabular}{lll}
\toprule
\textbf{Rango} & \textbf{Tipo} & \textbf{Uso} \\
\midrule
10.0.0.0/8           & Privada clase A  & Redes corporativas grandes \\
172.16.0.0/12        & Privada clase B  & Redes medianas \\
192.168.0.0/16       & Privada clase C  & Redes domésticas / pequeñas \\
127.0.0.0/8          & Loopback         & 127.0.0.1 = ``yo mismo'' \\
169.254.0.0/16       & APIPA link-local & Asignada automáticamente sin DHCP \\
100.64.0.0/10        & Shared Address   & NAT de operador (carrier-grade NAT) \\
0.0.0.0/0            & Default route    & ``todo el resto'' en tablas de ruteo \\
192.0.2.0/24         & Documentación    & Ejemplos en RFCs (RFC 5737) \\
\bottomrule
\end{tabular}
\caption{Rangos de IP especiales}
\end{table}

\section{\eh{right-arrow-curving-up} Routing --- Enrutamiento}

\begin{concepto}
Un \textbf{router} selecciona el mejor camino hacia el destino basándose
en su \textbf{tabla de ruteo}. Cada entrada en la tabla dice:
``para llegar a la red X, envía por la interfaz Y hacia el next-hop Z''.
La ruta más específica (mayor prefijo) siempre gana --- esto se llama
\textbf{Longest Prefix Match}.
\end{concepto}

\subsection{Routing estático}

\begin{codigo}
\begin{lstlisting}
! Cisco IOS - ruta estática
ip route 10.0.2.0 255.255.255.0 192.168.1.1
!            red destino   máscara      next-hop

! Ruta default (gateway de último recurso)
ip route 0.0.0.0 0.0.0.0 203.0.113.1

! Ver tabla de ruteo
show ip route
\end{lstlisting}
\end{codigo}

\subsection{OSPF --- Open Shortest Path First (RFC 2328)}

\begin{concepto}
\textbf{OSPF} es un protocolo de routing dinámico de \textbf{link-state}.
Cada router aprende la topología completa de la red y calcula el camino
más corto usando el \textbf{algoritmo de Dijkstra (SPF)}.

\emoji{light-bulb} \textbf{Analogía:} Cada router tiene un mapa completo de
la ciudad. Cuando necesita enviar un paquete, calcula la ruta más corta
por su cuenta usando el mapa. Si una calle se corta, actualiza su mapa
y recalcula. RIP en cambio es como preguntar a cada vecino cuántos pasos
hay a cada destino --- más simple pero menos preciso y lento en converger.
\end{concepto}

Conceptos clave de OSPF:
\begin{itemize}
    \item \textbf{Áreas:} la red se divide en áreas para reducir el tamaño de la LSDB.
          El \textbf{Área 0} (backbone) es obligatoria; todas las demás áreas
          deben conectarse a ella.
    \item \textbf{LSA} (Link State Advertisement): mensajes que intercambian los routers
          para describir sus enlaces y vecinos.
    \item \textbf{LSDB} (Link State Database): la base de datos de toda la topología.
          Todos los routers en el mismo área tienen la misma LSDB.
    \item \textbf{DR/BDR} (Designated Router / Backup DR): en redes multi-acceso
          (Ethernet), un solo router centraliza los LSAs para reducir tráfico.
    \item \textbf{Costo OSPF}: $\text{costo} = \frac{10^8}{\text{bandwidth en bps}}$.
          Ethernet 100 Mbps → costo 1. Serial 1.5 Mbps → costo 64.
\end{itemize}

\subsection{BGP --- Border Gateway Protocol (RFC 4271)}

\begin{concepto}
\textbf{BGP} es el protocolo de routing de Internet. Conecta
\textbf{Sistemas Autónomos (AS)} --- grandes redes de ISPs y corporaciones.
Es el único protocolo de routing que funciona a escala de todo Internet:
$>$900,000 prefijos IPv4 en la tabla BGP global.

\emoji{light-bulb} \textbf{Analogía:} Si OSPF es el GPS dentro de tu ciudad
(todos comparten el mismo mapa), BGP es el sistema de tratados entre países:
cada país (AS) decide sus propias políticas de tránsito y las anuncia
a sus vecinos mediante acuerdos (sessions BGP).
\end{concepto}

\begin{itemize}
    \item \textbf{eBGP:} entre routers de diferente AS (External BGP)
    \item \textbf{iBGP:} entre routers del mismo AS (Internal BGP)
    \item \textbf{AS number:} identificador del sistema autónomo (ej: AS13335 = Cloudflare)
\end{itemize}

\begin{cuidado}
\textbf{BGP Hijacking:} un AS puede anunciar prefijos IP que no le pertenecen,
robando o espiando tráfico. Caso famoso: Pakistan Telecom (2008) dejó YouTube
inaccesible \textit{en todo el mundo} durante 2 horas por un error de configuración BGP.

Defensa moderna: \textbf{RPKI} (Resource Public Key Infrastructure) ---
los dueños legítimos firman criptográficamente sus prefijos. Routers con
RPKI rechazan anuncios no firmados. Ver capítulo de seguridad.
\end{cuidado}

\subsection{HSRP --- Hot Standby Router Protocol}

\begin{concepto}
\textbf{HSRP} (RFC 2281, propietario Cisco) proporciona \textbf{gateway redundante}:
varios routers comparten una IP virtual. Si el router activo falla,
el standby toma la IP virtual en segundos. Los hosts solo ven la IP virtual.

Estándar abierto equivalente: \textbf{VRRP} (RFC 5798).
\end{concepto}

\section{\eh{recycle} NAT --- Network Address Translation}

\begin{concepto}
\textbf{NAT} permite que múltiples dispositivos con IPs privadas compartan
una sola IP pública. Es la razón por la que todos los dispositivos de tu
casa tienen IPs \texttt{192.168.x.x} pero salen a Internet con una sola IP.

\emoji{light-bulb} \textbf{Analogía:} El router es la recepción de una gran
empresa. Todas las llamadas internas salen con el número de la empresa, y la
recepcionista (NAT table) sabe a qué extensión interna redirigir cada respuesta.
\end{concepto}

\textbf{PAT} (Port Address Translation) = NAT overload: usa números de puerto
para distinguir entre múltiples conexiones internas que comparten la misma IP pública:

\begin{lstlisting}
  Interno                    Router NAT               Internet
  192.168.1.10:54321  →  203.0.113.1:10001  →   servidor:443
  192.168.1.11:52010  →  203.0.113.1:10002  →   servidor:80
  192.168.1.10:55000  →  203.0.113.1:10003  →   otro-server:443

  NAT Table:
  Entrada interna       ↔  Entrada pública
  192.168.1.10:54321   ↔  203.0.113.1:10001
  192.168.1.11:52010   ↔  203.0.113.1:10002
\end{lstlisting}

\section{\eh{globe-with-meridians} IPv6 (RFC 8200, 2017)}

\begin{concepto}
IPv4 tiene $\approx 4.3$ billones de direcciones y ya se agotaron (2011).
\textbf{IPv6} tiene 128 bits → $3.4 \times 10^{38}$ direcciones.
Suficientes para dar \textbf{670 cuatrillones de IPs por milímetro cuadrado}
de la superficie terrestre.
\end{concepto}

Formato: \texttt{2001:0db8:85a3:0000:0000:8a2e:0370:7334}

Reglas de abreviación:
\begin{itemize}
    \item Ceros iniciales en cada grupo se omiten: \texttt{0001} → \texttt{1}
    \item Grupos consecutivos de \texttt{0000} se reemplazan una vez por \texttt{::}
    \item \texttt{2001:0db8:0000:0000:0000:0000:0000:0001} → \texttt{2001:db8::1}
\end{itemize}

\begin{table}[h!]
\centering
\renewcommand{\arraystretch}{1.3}
\begin{tabular}{lll}
\toprule
\textbf{Tipo} & \textbf{Prefijo} & \textbf{Equivalente IPv4} \\
\midrule
Loopback        & \texttt{::1/128}     & 127.0.0.1 \\
Link-local      & \texttt{fe80::/10}   & 169.254.0.0/16 \\
Global unicast  & \texttt{2000::/3}    & IP pública \\
Unique local    & \texttt{fc00::/7}    & RFC 1918 (privadas) \\
Multicast       & \texttt{ff00::/8}    & 224.0.0.0/4 \\
\bottomrule
\end{tabular}
\caption{Tipos de direcciones IPv6}
\end{table}

\begin{cuidado}
IPv6 \textbf{no tiene broadcast}. Usa \textbf{multicast} para las funciones
que en IPv4 usaban broadcast. Por ejemplo:
\begin{itemize}
    \item ARP es reemplazado por \textbf{NDP} (Neighbor Discovery Protocol)
          usando mensajes multicast \texttt{ff02::1} (todos los nodos).
    \item Los routers escuchan en \texttt{ff02::2} (todos los routers).
\end{itemize}
\end{cuidado}

\subsection{ICMP --- Internet Control Message Protocol}

\begin{concepto}
\textbf{ICMP} (RFC 792) es el protocolo de mensajes de control de IP.
No transporta datos de usuario, sino mensajes de diagnóstico y error.
\end{concepto}

\begin{table}[h!]
\centering
\renewcommand{\arraystretch}{1.3}
\begin{tabular}{lll}
\toprule
\textbf{Tipo} & \textbf{Código} & \textbf{Mensaje} \\
\midrule
8  & 0 & Echo Request (ping enviado) \\
0  & 0 & Echo Reply (respuesta a ping) \\
11 & 0 & Time Exceeded (TTL=0, usado por traceroute) \\
3  & 0 & Destination Unreachable - Network unreachable \\
3  & 1 & Destination Unreachable - Host unreachable \\
3  & 3 & Destination Unreachable - Port unreachable \\
\bottomrule
\end{tabular}
\caption{Mensajes ICMP comunes}
\end{table}

\begin{ejemplo}
\textbf{¿Cómo funciona traceroute?}

Traceroute envía paquetes UDP (o ICMP en Windows) con TTL=1, 2, 3...
\begin{itemize}
    \item TTL=1: el primer router lo decrementa a 0 y responde con ICMP
          ``Time Exceeded'' → revela la IP del primer salto.
    \item TTL=2: el segundo router responde → IP del segundo salto.
    \item Y así sucesivamente hasta llegar al destino, que responde con
          ICMP ``Port Unreachable'' (o Echo Reply en ICMP mode).
\end{itemize}
\end{ejemplo}
